\label{chap:83-07}

The connection is constructed from the complete orbital classification and from
its first coinductive elevation. The following development fixes the three
initial histories, the oriented boundaries and the chain
\(w_6\to w_{30}\to R_{36}\) before introducing the nonadic holonomy.

\begin{HMTScientificOwnerSubordinated}
\input{ampliaciones_sucesoras_20260824/parte_i_ii/owners/U008_orbitas_elevacion_r36}
\end{HMTScientificOwnerSubordinated}

% Unidad U009. Conexión nonádica conjunta; redacción autónoma.

\section{Joint nonadic connection, compression, and terminal memory}
\label{p12-83-c7-s1:chap:conexion-nonadica-conjunta}

The nine phases that appear after \(R_{36}\) are not nine independent
proofs. They are the nine local components of a single connection
on enriched states. Their composition around a turn defines
a holonomy. The return of the phase label conserves the continuity of
prefix, carry, boundary, Witt shadow, and memory.

\subsection{Transport of the enriched state by the nonadic holonomy and memory update}

At the level \(K\), we write
\[
 \widehat x_K=
 (B_K,C_K,p_K,c_K,\Sigma_K,W_K,g(K),q_{{\rm mem},K}),
\]
where:
\begin{enumerate}
 \item \(B_K=(u_K^\pi,u_K^e,u_K^\varphi)\) is the visible block of three
       words of six trits;
 \item \(C_K\) is the cylinder ternario compatible;
 \item \(p_K\) is the complete prefix;
 \item \(c_K\) is the accumulated carry;
 \item \(\Sigma_K\) is the signature joint of boundary;
 \item \(W_K\) is the Witt incidence shadow;
 \item \(g(K)=1+((K-1)\bmod9)\) is the public phase label;
 \item \(q_{{\rm mem},K}\in\mathbb Z_{\geq0}\) is the memory vertical not
       acotada.
\end{enumerate}
The dodecaphasic coordinate is derived:
\[
 \beta_K=[q_{{\rm mem},K}]_{12}.
\]
\(q_{\rm mem}\) is not replaced by \(\beta\), since reduction modulo
twelve would lose the number of circuits.

\subsection{Local relations and holonomy}

For each \(K\), let
\[
 \mathcal R_K\subseteq
 \widehat X_K\times\widehat X_{K+1}
\]
the relation of compatible extension between complete states. The relational composition of a loop is
\begin{equation}
 \Gamma_{9,K}
 :=\mathcal R_{K+8}\circ\cdots\circ\mathcal R_{K+1}\circ\mathcal R_K
 :\widehat X_K\longrightarrow\widehat X_{K+9}.
\label{p12-83-c7-s1:eq:gamma9-composicion-local}
\end{equation}

\begin{definition}[Nonadic holonomy]
\label{p12-83-c7-s1:def:holonomia-nonadica}
The holonomy of one turn is
\[
 \Gamma_9:=\operatorname{Hol}_{C_9}(\nabla_{\rm TPK}),
\]
whose realization at level \(K\) is
\(\Gamma_{9,K}\). The nine relations \(\mathcal R_K,\ldots,
\mathcal R_{K+8}\) are charts of this single connection.
\end{definition}

\begin{theorem}[Phase return without restart]
\label{p12-83-c7-s1:thm:retorno-nonadico-sin-reinicio}
If \((\widehat x_K,\widehat x_{K+9})\in\Gamma_{9,K}\), then
\[
 g(K+9)=g(K),
 \qquad
 q_{{\rm mem},K+9}=q_{{\rm mem},K}+1,
\]
and
\[
 \beta_{K+9}=\beta_K+1\pmod{12}.
\]
In general, \(\widehat x_{K+9}\neq\widehat x_K\).
\end{theorem}

\begin{proof}
The first identity follows from the definition of \(g\). A complete turn increments the unbounded winding counter by definition. Reducing that identity modulo twelve yields the third. The remaining coordinates are transported by the composition \cref{p12-83-c7-s1:eq:gamma9-composicion-local}; no law of the TPK resets them. Therefore, equality of phase does not imply equality of the complete state.
\end{proof}

\subsection{Exhaustive classification of the \(9\) phases of the Topological Phase Kernel}

The visible blocks and the local censuses of the first round are
\[
\begin{array}{c|rrrrr|ccc}
K&N_\pi&N_e&N_\varphi&N_{\rm loc}&N_{\rm ext}
 &u^\pi&u^e&u^\varphi\\\hline
1&378&422&349&3&2&012222&121221&112202\\
2&238&368&388&2&5&010211&102011&121020\\
3&184&284&173&5&4&002111&012222&010210\\
4&207&207&122&4&1&110221&102011&010200\\
5&39&151&151&1&1&222220&021222&102011\\
6&110&110&69&1&3&111201&201202&221021\\
7&81&29&81&3&3&212121&222210&200221\\
8&59&59&59&3&2&200121&212212&110110\\
9&43&43&43&2&1&100100&020112&010122
\end{array}.
\]

\begin{theorem}[Functions of the nine components]
\label{p12-83-c7-s1:thm:funciones-nueve-componentes}
In the preceding turn, the nine components perform the following
functions:
\begin{enumerate}
 \item local opening of three orientations and first resolution;
 \item binary split with a local maximum of five extensions;
 \item maximal local multiplicity of five states;
 \item contraction of four local states to an extension;
 \item first strong saturation, with a single continuation;
 \item unique local present and triple opening of the future;
 \item triple mirror fiber;
 \item synchronization of the three censuses in \(59\);
 \item oriented closure: two local input states and a relation that,
       after the additional horizon, conserves a prolongation.
\end{enumerate}
\end{theorem}

\begin{proof}
Each assertion is read from the censuses \(N_{\rm loc},N_{\rm ext}\), from the
equality or inequality of \(N_\pi,N_e,N_\varphi\), and from the explicit extension
relation in the ninth, tenth, and eleventh phases to be proved
in the following unit. The table enumerates the entire revolution; no subset
of phases has been chosen.
\end{proof}

The fifth component coincides with the first common boundary.
\[
 R_{36}=(222220|021222|102011)^T.
\]
The ninth component contains
\[
 B_9=(100100|020112|010122).
\]
The presence of two local states in the ninth phase and of three extensions
from the oriented state belongs to different layers; they are not contradictory
censuses.

Holonomy acts on the complete enriched state. Before its modal
quotient is published, one therefore constructs the involution that exchanges
the \(\pi\) and \(e\) fibres, the fixed \(\varphi\) axis, the ternary
invariants, and the local \(3\to1\) resolution determined by the second
horizon.

\begin{HMTScientificOwnerSubordinated}
\input{ampliaciones_sucesoras_20260824/parte_i_ii/owners/U010_espejo_seleccion_tres_a_uno}
\end{HMTScientificOwnerSubordinated}

\subsection{Joint transfer of the fiber of 468 emissions}

The factorization of the catalog is
\[
 \mathcal U_6^{\rm cat}\cong
 \mathfrak S_+\times\mathfrak S_\times,
 \qquad
 |\mathfrak S_+|=18,\qquad |\mathfrak S_\times|=26.
\]
For a function \(f\) on the catalogue, let
\[
 (E_+f)(s,e)=\frac1{18}\sum_{s'\in\mathfrak S_+}f(s',e),
\]
and
\[
 (E_\times f)(s,e)=\frac1{26}
 \sum_{e'\in\mathfrak S_\times}f(s,e').
\]
The projectors \(E_+\) and \(E_\times\) commute.

Let
\[
 \tau=(+1,-1,0,+1,-1,0,+1,-1,0),
 \qquad
 P_{+1}=E_+,quad P_{-1}=E_\times,quad P_0=I.
\]
On \(\mathcal U_6^{\rm cat}\times C_9\) we define phase transfer
by
\begin{equation}
 \mathcal K_{\rm ph}\bigl((u,r),(v,[r+1]_9)\bigr)
 :=P_{\tau(r)}(u,v),
 \label{p12-83-c7-s1:eq:u009-kph-explicita}
\end{equation}
and we make it zero for the remaining pairs of phases. Thus, the cyclic permutation
and the operator applied in each of the nine steps are determined.

\begin{proposition}[Joint Renewal After a Turn]
\label{p12-83-c7-s1:prop:transferencia-conjunta-nueve}
\[
 \mathcal K_{\rm ph}^{\,9}=E_+E_\times.
\]
In particular, on a point mass,
\[
 (E_+E_\times)(u,v)=\frac1{18\cdot26}=\frac1{468}.
\]
\end{proposition}

\begin{proof}
The ordered product over a loop is
\[
 P_{\tau(8)}P_{\tau(7)}\cdots P_{\tau(0)}
 =I E_\times E_+ I E_\times E_+ I E_\times E_+
 =E_+^3E_\times^3I^3=E_+E_\times.
\]
The second equality uses commutativity and the final one idempotence. The
second result is the uniform distribution on the product of the two
signatures.
\end{proof}

This identity expresses a joint renewal, not nine independent filtrations.

\subsection{Decomposition of the holonomy \(U_{\Gamma_9}\) into visible compression \(T\) and orthogonal component \(\eta\): \(I-T^*T=\eta^*\eta\)}

Let \(\mathcal H_{\rm vis}\) and \(\mathcal H_{\rm full}\) be Hilbert spaces for the visible and complete realizations. Let
\[
 J:\mathcal H_{\rm vis}\longrightarrow\mathcal H_{\rm full}
\]
an isometric inclusion and
\(U_{\Gamma_9}\) an isometric realization of the holonomy. We define
\[
 T:=J^*U_{\Gamma_9}J,
 \qquad
 \eta:=(I-JJ^*)U_{\Gamma_9}J.
\]

\begin{theorem}[Compression-decomposition expression]
\label{p12-83-c7-s1:thm:conexion-compresion-expresion}
The following hold
\[
 \boxed{U_{\Gamma_9}J=JT+\eta,}
\]
\[
 J^*\eta=0,
 \qquad
 \boxed{I-T^*T=\eta^*\eta.}
\]
\end{theorem}

\begin{proof}
We decompose \(U_{\Gamma_9}J\) by means of the orthogonal projections
\(JJ^*\) and \(I-JJ^*\):
\[
 U_{\Gamma_9}J
 =JJ^*U_{\Gamma_9}J+(I-JJ^*)U_{\Gamma_9}J
 =JT+\eta.
\]
Orthogonality gives \(J^*\eta=0\). Since \(J\) and
\(U_{\Gamma_9}\) are isometries,
\[
 I=J^*U_{\Gamma_9}^*U_{\Gamma_9}J
 =T^*T+\eta^*\eta.
\]
Reorder the proof of the last identity.
\end{proof}

\(T\) is the visible compression; \(\eta\) is the complementary expression.
Neither should be called a strict contraction. For that function a
different operator is reserved.

\subsection{Strict contraction and terminal telescoping}

Let \(V_9\) be an isometry and let \(q\) be a parameter with \(0<q<1\). We define
\[
 M_9=qV_9,
 \qquad
 D_9=(I-M_9^*M_9)^{1/2}.
\]
Then \(\|M_9\|=q<1\) and
\[
 I-M_9^*M_9=D_9^*D_9.
\]

\begin{theorem}[Stein Identity with Terminal Term]
\label{p12-83-c7-s1:thm:stein-telescopia-terminal}
For every \(L\geq1\),
\[
 \boxed{
 I=
 \sum_{r=0}^{L-1}M_9^{*r}D_9^*D_9M_9^r
 +M_9^{*L}M_9^L.}
\]
More generally, if \(S_0\) satisfies
\(S_0-M_9^*S_0M_9=D_9^*D_9\), then
\[
 \boxed{
 S_0=
 \sum_{r=0}^{L-1}M_9^{*r}D_9^*D_9M_9^r
 +M_9^{*L}S_0M_9^L.}
\]
\end{theorem}

\begin{proof}
We multiply the defect identity on the left by \(M_9^{*r}\) and on the
right by \(M_9^r\). Upon summing over \(r=0,\ldots,L-1\), the interior terms
telescope, leaving the initial and terminal terms. The same proof
applies to the general Stein equation.
\end{proof}

The final term does not vanish at finite depth. Its disappearance in a
limit requires a convergence proof; it is not part of the algebraic
identity.

\subsection{Carry Cocycle and Relational Powers}

For composable arrows \(\gamma_1,\gamma_2\), the carry satisfies
\[
 \operatorname{Carr}(\gamma_2\gamma_1)
 =\operatorname{Carr}(\gamma_2)H(\gamma_1)
 +Y(\gamma_2)\operatorname{Carr}(\gamma_1).
\]
The formula preserves the contribution of the first stage after transporting it through the second.

The longer loops are written with displaced indices:
\begin{align*}
 \Gamma_{27,K}
 &=\Gamma_{9,K+18}\circ\Gamma_{9,K+9}\circ\Gamma_{9,K},\\
 \Gamma_{54,K}
 &=\Gamma_{9,K+45}\circ\cdots\circ\Gamma_{9,K},\\
 \Gamma_{108,K}
 &=\Gamma_{9,K+99}\circ\cdots\circ\Gamma_{9,K}.
\end{align*}
Only after declaring a global operator on the graded union may
this notation be abbreviated as \(\Gamma_9^3,\Gamma_9^6,\Gamma_9^{12}\).

\subsection{Complete nonadic state: phase, carry, boundary, route, and memory after the prolongation}

The one-turn connection is recorded by the quintuple
\[
 \mathfrak N_9=
 \bigl(\Gamma_9,T,\eta,\operatorname{Carr},S_{\rm term}\bigr),
\]
where \(S_{\rm term}=M_9^{*L}S_0M_9^L\) at a finite depth
\(L\). An application functor that uses the connection must declare its
compatibility with all five fields. Commuting only with the visible phase does not
constitute naturality with respect to the complete nonadic state.

\subsection{Obstructions of the holonomy \(\Gamma_9\): return \(9\to1\), composition indices and preservation of phase, carry, boundary, route and memory}

\begin{criterion}[Refutation Criteria of the Connection]
The construction is refuted if:
\begin{enumerate}
 \item the nine phases are presented as independent filters;
 \item the return \(9\to1\) resets \(q_{\rm mem}\), the prefix, or the carry;
 \item \(T\) is called a strict contraction;
 \item some identity of
       \cref{p12-83-c7-s1:thm:conexion-compresion-expresion} fails;
 \item the terminal term is removed at finite depth;
 \item \(\Gamma_{27,K}\) is composed without shifting the indices;
 \item an application functor ignores any of the five fields of
       \(\mathfrak N_9\).
\end{enumerate}
\end{criterion}

The complete developments that follow first reconstruct the finite emitter,
its factorization and its microfibers; they then unfold the specular involution,
the three extensions of \(B_9\), uniqueness at two horizons and the
Bockstein--Hensel inverse that preserves memory during phase return.

\begin{HMTScientificOwnerSubordinated}
\input{sections/hmt/03b_alfabeto_factorizacion}
\end{HMTScientificOwnerSubordinated}

\begin{HMTScientificOwnerSubordinated}
\input{sections/hmt/14_numero_heisenberg_y_d108}
\end{HMTScientificOwnerSubordinated}
